The \xmlNode{uncertainty} node makes a value a random variable. It holds exactly one distribution
node, chosen from the distributions listed below. HERON samples a new value from the distribution
for each realization of the outer workflow, so its effect appears in the statistics HERON reports
over those realizations. Using \xmlNode{uncertainty} in a cash flow forces HERON to use the bilevel
(outer and inner) workflow.

HERON currently samples \xmlNode{uncertainty} only for these \xmlNode{CashFlow} parameters:
\xmlNode{driver}, \xmlNode{reference\_price}, \xmlNode{reference\_driver}, and
\xmlNode{scaling\_factor\_x}. Other nodes accept \xmlNode{uncertainty} as input, but it is not
sampled there and should not be used.

For example, a reference price drawn from a normal distribution:
\begin{lstlisting}[style=XML]
  <reference_price>
    <uncertainty>
      <Normal name="price_dist">
        <mean>10</mean>
        <sigma>0.05</sigma>
      </Normal>
    </uncertainty>
  </reference_price>
\end{lstlisting}
The \xmlAttr{name} of the distribution must be given but is replaced by HERON when it writes the
RAVEN workflow.

Every distribution accepts the optional \xmlNode{lowerBound} and \xmlNode{upperBound} subnodes,
which truncate the distribution. Parameters marked ``Required'' must be given. Distributions not
listed here are deprecated: HERON passes them to RAVEN without checking them and warns that they
will be removed.
